\subsection{Controlling the Time Step Size in GFR}
%\subsection{\texorpdfstring%
%            {\colorbox{yellow}{Controlling the Time Step Size in GFR}}%
%                              {Controlling the Time Step Size in GFR}}
\label{SEC:Time:Controlling}

%%%%%%%%%%%%%%%%%%%%%%%%%%%%%%%%%%%%%%%%%%%%%%%%%%%%%%%%%%%%%%%%%%%%%%%%%%%%%%%%

%\subsubsection{\texorpdfstring{\colorbox{yellow}{Ramping}}{Ramping}}
\subsubsection{Ramping}
\label{SEC:Time:Controlling:Ramping}

%%%%%%%%%%%%%%%%%%%%%%%%%%%%%%%%%%%%%%%%%%%%%%%%%%%%%%%%%%%%%%%%%%%%%%%%%%%%%%%%

After solution initialization, the flow may have very strong gradients
as the boundary conditions are applied.
%
This may necessitate a small timestep to keep the solution stable
during these initial transients.
%
However, as the flow progresses, a larger timestep may be permissible.

GFR has a ramping capability that applies a multiplicative factor to
either the requested constant CFL or constant timestep.
%
This factor linearly varies between specified beginning and ending
values over a certain period of iterations.

\SeeAlso{
CFL\_BEG,
CFL\_END,
CFL\_CYCLES
}

%%%%%%%%%%%%%%%%%%%%%%%%%%%%%%%%%%%%%%%%%%%%%%%%%%%%%%%%%%%%%%%%%%%%%%%%%%%%%%%%

\subsubsection{Factors that Indirectly Affect Maximum Stable Time Step Size}%
\label{SEC:Time:Controlling:Indirectly}

The following are some considerations to be mindful of with how they affect the
maximum stable time step.

%%%%%%%%%%%%%%%%%%%%%%%%%%%%%%%%%%%%%%%%%%%%%%%%%%%%%%%%%%%%%%%%%%%%%%%%%%%%%%%

\subsubsubsection{Effect of Limiting\slash Filtering on Time Step Size}
\label{SEC:Time:Controlling:Indirectly:Limiting}

Limiting\slash filtering will generally increase the maximum stable time step
size, since both methods essentially damp the higher order polynomial modes.
%
This decreases the solution accuracy, but permits the solution to advance
further in time without those higher order modes amplifying into a numerical 
instability.
%
Note, however, that a larger time step size will likely cause some additional
time-step induced instabilities, which may cause more cells to be marked and
limited\slash filtered to maintain stability.
%
Selecting the optimal combination of limiter/filter and time step settings
can be a tedious process.

%%%%%%%%%%%%%%%%%%%%%%%%%%%%%%%%%%%%%%%%%%%%%%%%%%%%%%%%%%%%%%%%%%%%%%%%%%%%%%%

\subsubsubsection{Effect of Solution Order on Time Step Size}
\label{SEC:Time:Controlling:Indirectly:P_vs_dt}

For high-order flux reconstruction, discontinuous Galerkin, and spectral-element
methods, it has been shown that the maximum stable time step scales with the
solution polynomial order as $1/\p{}^2$\cite{Hesthaven_2010, Huynh_2014}.
%
Thus, for higher-order solutions, expect substantially smaller time step sizes.
%
Further details including a rigorous mathematical discussion can be found in
Section 4.7 of \linecite{Hesthaven_2010}.

%%%%%%%%%%%%%%%%%%%%%%%%%%%%%%%%%%%%%%%%%%%%%%%%%%%%%%%%%%%%%%%%%%%%%%%%%%%%%%%

\subsubsubsection{Effect of Solution Order on CFL}
\label{SEC:Time:Controlling:Indirectly:P_vs_CFL}

This is most likely a GFR-specific behavior that may be related to how the CFL
is computed, but the solution order affects the maximum stable CFL inversely to
how it affects the maximum stable time step size.
%
That is for higher-order solutions, larger CFL values may be possible.
%
As an example, take the spacecraft cabin ventilation fan shown in
\fref{FIG:Files:PW:Spacefan}.
%
\Tref{TBL:Time:Spacefan} shows the maximum CFL used, and resulting physical time
step size, for a single configuration of the fan run at \PO[1]--\PO[4], \thatis
the same grid and flow conditions were used and the only thing different between
the simulations was the solution order.
%
All of these simulations used the 11-stage\slash\nth{3}-order VOERK scheme
(\thatis \inp[nmlTimestep::RungeKuttaScheme]{Runge\_Kutta\_Scheme}=113).
%
\begin{table}[tb]
  \centering
  \caption[Example maximum stable CFL relative to solution order.]%
          {Maximum stable \inptext{CFL} and resulting physical time step size,
           $\Delta t$, for simulations of the spacecraft cabin ventilation fan
           relative to the solution order, $\PO$.}
  \label{TBL:Time:Spacefan}
  \begin{tabular}{ @{} @{\quad} c @{\qquad} c @{\qquad} c @{\quad} @{} }
    \PO & \inptext{CFL} & \dt{} [s] \\
    \midrule
    1 & 1.5 & \num{9.2e-9} \\
    2 & 3.5 & \num{6.0e-9} \\
    3 & 5.5 & \num{4.3e-9} \\
    4 & 8.0 & \num{3.1e-9}
  \end{tabular}
\end{table}
%
This example clearly shows how the maximum \inptext{CFL} that keeps GFR stable
increases with \PO, but the corresponding time step size still decreases with
increasing \PO{} due to the $1/\p{}^2$ scaling from
\sref{SEC:Time:Controlling:Indirectly:P_vs_dt}.

%%%%%%%%%%%%%%%%%%%%%%%%%%%%%%%%%%%%%%%%%%%%%%%%%%%%%%%%%%%%%%%%%%%%%%%%%%%%%%%

\subsubsubsection{Effect of Transients on Time Step Size}
\label{SEC:Time:Controlling:Indirectly:Transients}

When starting a new simulation or restarting a simulation with significant
changes to the settings (\forexample changing the solution order or adjusting
the pressure for a boundary condition), transients are created in the flow field
that generally contain larger and sharper gradients than would normally exist in
the well-established, stable flow field.
%
These disturbances may require the use of a smaller time step size to prevent
them from growing into numerical instabilities.
%
Once the disturbances have subsided or convected out of the domain, it may
be possible to increase the time step size.

%%%%%%%%%%%%%%%%%%%%%%%%%%%%%%%%%%%%%%%%%%%%%%%%%%%%%%%%%%%%%%%%%%%%%%%%%%%%%%%%

\subsubsection{Finding a Near-Optimal Stable CFL}%
\label{SEC:Time:Controlling:OptimizingCFL}

%%%%%%%%%%%%%%%%%%%%%%%%%%%%%%%%%%%%%%%%%%%%%%%%%%%%%%%%%%%%%%%%%%%%%%%%%%%%%%%%

This section details one method for trying to find a maximum stable
CFL\slash time step size for a given simulation.

%%%%%%%%%%%%%%%%%%%%%%%%%%%%%%%%%%%%%%%%%%%%%%%%%%%%%%%%%%%%%%%%%%%%%%%%%%%%%%%%

After creating the grid, the first step of running a simulation is creating the
input file and running some initial tests to verify there are no issues with the
grid or input file.
%
Perform tests with a smaller CFL value that is almost certain to be stable, say
0.1-0.5, and have them run for 100 to 10,000 time steps.

%%%%%%%%%%%%%%%%%%%%%%%%%%%%%%%%%%%%%%%%%%%%%%%%%%%%%%%%%%%%%%%%%%%%%%%%%%%%%%%%

Once it seems that there are no major issues, this is when to begin the second
step of trying to run the simulation for a large number of time steps.
%
During this phase, initial transients will dominate the solution and this is the
time to play with the filter\slash limiting settings in order to find a
combination that keeps the simulation stable without damping out too much of the
smallest resolved scales.

%%%%%%%%%%%%%%%%%%%%%%%%%%%%%%%%%%%%%%%%%%%%%%%%%%%%%%%%%%%%%%%%%%%%%%%%%%%%%%%%

After the initial transients in the simulation are mostly eliminated and the
solution appears to be stable, the next step is to push the CFL in order to find
the maximum stable time step.
%
If the simulation is running, the interrupt file \file{ADJUST\_CFL}
(\sref{SEC:Files:Interrupt:Adjust_CFL}) can be used to modify the CFL input
variables; otherwise, the CFL inputs can simply be adjusted in the input file
for the next simulation.
%
Utilize the CFL ramping inputs to slowly increase the CFL with the actual goal
of causing the simulation to diverge and crash.
%
If the previous simulation was using a CFL of 0.5, one could set up the next
simulation with the inputs
%
\begin{InputExample}
  \item[{\inptext{CFL\_Ramping\_Beg = 0.5}}]
  \item[{\inptext{CFL\_Ramping\_End = 15.0}}]
  \item[{\inptext{CFL\_Ramping\_Iterations = 14500}}]
\end{InputExample}
%
The final CFL value of 15.0 is chosen because it is assumed that this would
almost certainly be too large for the simulation to remain stable.
%
This new simulation is run and say it crashes once it reaches a CFL of 10.2.
%
One should generally reduce the final CFL number by 10--15\%, and use this
value as the new CFL going forward.
%
If you want to push the CFL to the highest possible value, you can also search
through the output file \file{residual.dat} in order to find the approximate
iteration at which the simulation residual begins to diverge.
%
However, it is generally less of a headache to trim a few extra percent off
just to make sure the simulation remains stable. 

%%%%%%%%%%%%%%%%%%%%%%%%%%%%%%%%%%%%%%%%%%%%%%%%%%%%%%%%%%%%%%%%%%%%%%%%%%%%%%%%

In this example, the number of CFL ramping iterations is chosen to be 14,500 so
that each integer increment of the CFL number takes 1,000 iterations.
%
This offers a moderate level of resolution in terms of what CFL value will
maintain stability.
%
For example, if the CFL ramping is over 1,000 iterations, narrowing down the
optimum CFL value would be less accurate.

%%%%%%%%%%%%%%%%%%%%%%%%%%%%%%%%%%%%%%%%%%%%%%%%%%%%%%%%%%%%%%%%%%%%%%%%%%%%%%%%

It is also important to remember that the maximum stable time step is something
that can change over the course of a simulation since it is affected by many
different aspects.
%
If a simulation appears to be stable at a given CFL value and then suddenly
crashes, try restarting from a recent restart file (but not too recent in case
the instability already exists in the restart file) with a smaller CFL (say 50\%
smaller).
%
If the simulation crashes at nearly the same simulation time, you might want
to try an even smaller time step with an even earlier restart file, but most
likely the instability is due to something that will need to be treated with
numerical methods like de-aliasing or stronger filtering\slash limiting.

%%%%%%%%%%%%%%%%%%%%%%%%%%%%%%%%%%%%%%%%%%%%%%%%%%%%%%%%%%%%%%%%%%%%%%%%%%%%%%%%

With enough practice and experience (and patience), various aspects of the steps
provided in this section can get blended together to speed up this process even
more.
%
Unfortunately, a good understanding and feel for how the time step size responds
to the factors listed in \sref{SEC:Time:Controlling:Indirectly} is really needed
to blend these steps effectively.

%%%%%%%%%%%%%%%%%%%%%%%%%%%%%%%%%%%%%%%%%%%%%%%%%%%%%%%%%%%%%%%%%%%%%%%%%%%%%%%%

